\documentclass[10pt,a4paper]{article}

\usepackage[T1]{fontenc}
\usepackage[utf8]{inputenc}
\usepackage[ngerman]{babel}
\usepackage[a4paper,margin=14mm]{geometry}
\usepackage{lmodern}
\usepackage[table]{xcolor}
\usepackage{tabularx}
\usepackage{array}
\usepackage{microtype}
\usepackage[most]{tcolorbox}

\pagestyle{empty}
\setlength{\parindent}{0pt}
\setlength{\parskip}{0pt}
\renewcommand{\familydefault}{\sfdefault}

\definecolor{Navy}{HTML}{17324D}
\definecolor{Blue}{HTML}{3B6EA5}
\definecolor{Sky}{HTML}{EAF3FA}
\definecolor{Mint}{HTML}{E8F5EF}
\definecolor{Gold}{HTML}{F3B33D}
\definecolor{Ink}{HTML}{1E2935}
\definecolor{Muted}{HTML}{617080}
\definecolor{Line}{HTML}{D8E1E8}
\definecolor{CodeBg}{HTML}{F5F7F9}

\color{Ink}

\newcommand{\sectiontitle}[1]{%
  \vspace{2.3mm}
  {\color{Blue}\bfseries\large #1}\par
  \vspace{0.8mm}
  {\color{Blue}\rule{\linewidth}{0.7pt}}\par
  \vspace{1.2mm}
}

\newcommand{\classcard}[2]{%
  \fcolorbox{Blue}{white}{%
    \parbox{63mm}{%
      \centering\vspace{1mm}{\large\bfseries\texttt{#1}}\par
      \vspace{1mm}{\color{Blue}\rule{\linewidth}{0.6pt}}\par
      \vspace{1mm}\raggedright\ttfamily\small #2\vspace{1mm}
    }%
  }%
}

\begin{document}

\colorbox{Navy}{%
  \parbox{\dimexpr\linewidth-2\fboxsep\relax}{%
    \vspace{3mm}\color{white}
    \begin{tabularx}{\linewidth}{@{}Xr@{}}
      {\small\bfseries INFORMATIK · KLASSE 10} &
      \colorbox{Gold}{\color{Navy}\small\bfseries\strut LÖSUNGEN}
    \end{tabularx}
    \vspace{2.5mm}

    {\fontsize{22}{25}\selectfont\bfseries Sicherungsblatt: Aufgabe 2}\par
    \vspace{1mm}
    {\large Redundanzen in Datenbanken}
    \vspace{3mm}
  }%
}

\vspace{3mm}
\colorbox{Sky}{%
  \parbox{\dimexpr\linewidth-2\fboxsep\relax}{%
    \vspace{1.8mm}
    \textcolor{Blue}{\bfseries Ziel:}
    Du erkennst redundante Daten, erklärst ihre Probleme und kannst Daten
    sinnvoll auf mehrere Tabellen aufteilen.
    \vspace{1.8mm}
  }%
}

\sectiontitle{1 · Redundanz erkennen}

\textbf{Redundanz} bedeutet: \textbf{Dieselbe Information wird mehrfach gespeichert.}

\vspace{1.5mm}
\renewcommand{\arraystretch}{1.2}
\begin{tabularx}{\linewidth}{@{}>{\ttfamily\small}p{16mm}>{\ttfamily\small}p{25mm}>{\ttfamily\small}p{47mm}X@{}}
  \rowcolor{Sky}\bfseries photo\_id & \bfseries username & \bfseries email & \bfseries description \\
  41 & mia & mia@example.org & Sonnenuntergang \\
  \rowcolor{CodeBg}42 & mia & mia@example.org & Mein Fahrrad \\
  43 & mia & mia@example.org & Ausflug am See \\
\end{tabularx}

\vspace{1.4mm}
\colorbox{CodeBg}{\parbox{\dimexpr\linewidth-2\fboxsep\relax}{%
  \texttt{username} und \texttt{email} werden für jedes Foto wiederholt gespeichert.
  Das ist redundant.
}}

\sectiontitle{2 · Warum sind redundante Daten problematisch?}

\begin{tabularx}{\linewidth}{@{}p{7mm}X p{7mm}X@{}}
  \textcolor{Gold}{\Large\bfseries !} & Daten werden unnötig mehrfach gespeichert. &
  \textcolor{Gold}{\Large\bfseries !} & Speicher- und Pflegeaufwand steigen. \\[1mm]
  \textcolor{Gold}{\Large\bfseries !} & Änderungen müssen an mehreren Stellen erfolgen. &
  \textcolor{Gold}{\Large\bfseries !} & Vergessene Änderungen führen zu widersprüchlichen, inkonsistenten Daten. \\
\end{tabularx}

\vspace{1.5mm}
\colorbox{Sky}{\parbox{\dimexpr\linewidth-2\fboxsep\relax}{%
  \textcolor{Blue}{\bfseries Beispiel:}
  Ändert Mia ihre E-Mail-Adresse nur bei einem Foto, stehen anschließend für
  dieselbe Benutzerin verschiedene E-Mail-Adressen in der Tabelle.
}}

\sectiontitle{3 · Bessere Lösung: Daten auf zwei Tabellen aufteilen}

\begin{center}
  \classcard{users}{id\par username\par email}
  \hspace{18mm}
  \classcard{photos}{id\par user\_id[users]\par description\par url\par created\_at\par updated\_at}
\end{center}

\vspace{1mm}
\begin{tabularx}{\linewidth}{@{}X X@{}}
  \textcolor{Blue}{\bfseries \texttt{users}:} Benutzerinformationen werden hier einmal gespeichert. &
  \textcolor{Blue}{\bfseries \texttt{photos}:} Jedes Foto enthält nur seine Fotoinformationen und die passende Zuordnung. \\
\end{tabularx}

\sectiontitle{4 · Merke}

\begin{tcolorbox}[colback=white,colframe=Line,arc=3mm,boxrule=0.7pt,left=3mm,right=3mm,top=2mm,bottom=2mm]
  - Daten von Fotos werden in einer separaten Tabelle gespeichert, um
  Redundanzen innerhalb der Tabelle \texttt{users} zu vermeiden.\par
  - Jeder Sachverhalt wird möglichst nur einmal gespeichert.\par
  - Benutzerinformationen gehören in \texttt{users}, Fotoinformationen in \texttt{photos}.\par
  - Über \texttt{users.id} und \texttt{photos.user\_id} lassen sich passende Datensätze wieder zuordnen.
\end{tcolorbox}

\vfill
{\color{Line}\rule{\linewidth}{0.5pt}}\par
\vspace{1.5mm}
\begin{tabularx}{\linewidth}{@{}Xr@{}}
  {\small\color{Muted}Klasse 10 · Datenbanken} &
  {\small\color{Muted}Sicherungsblatt · Aufgabe 2}
\end{tabularx}

\end{document}
